\chapter*{Abstract}

This thesis proposes an interpretable method for data-to-text generation from structured data. The generation process is divided into two stages: constructing a meaning representation and performing rule-based surface realization. This enables automatic or human evaluation of text semantics and explicit tracing of the text-generation process, rather than relying on a single, non-interpretable generation step.

Three pipelines for constructing meaning representations as RDF triples using large language models are proposed and evaluated. The Text-First Extraction with Normalization pipeline achieves the strongest reported results, while the remaining methods provide greater control over the predicate set and auditable intermediate artifacts. A domain-specific model is also trained to distill the operation of the complete processing pipelines into a single inference step. Training with QLoRA adapters substantially improves triple extraction over the base model, although performance depends on the domain and evaluation metric.

For surface realization, the thesis proposes using OpenEvolve to construct verbalization rules expressed as executable Python code. The evolutionary process uses LLM-generated code mutations and automatic text-evaluation metrics as objective functions. Experiments conducted on the WebNLG dataset show that optimizing reference-based evaluation metrics and metrics based on the LLM-as-Judge paradigm produces surface realizers with different properties. The evolved programs do not surpass the text quality achieved by a fine-tuned BART model or by zero-shot LLM generation, but they provide transparent and reproducible realization logic that can be inspected and modified without using an LLM at inference time.

The results confirm that explicit intermediate representations and executable generation artifacts are practical tools for interpretable data-to-text generation.
